\documentclass[conference]{IEEEtran}

\usepackage{iftex}
\ifPDFTeX
  \usepackage[utf8]{inputenc}
  \usepackage[T1]{fontenc}
\fi
\usepackage{amsmath,amssymb,bm,mathtools}
\usepackage{newtxtext,newtxmath}
\usepackage{graphicx,svg,booktabs,adjustbox}
\usepackage{siunitx}
\usepackage[final]{microtype}
\usepackage{xcolor}
\usepackage{cite}
\usepackage[breaklinks=true,hidelinks]{hyperref}
\usepackage[nameinlink,capitalise,noabbrev]{cleveref}

\interdisplaylinepenalty=2500
\sisetup{detect-weight=true,detect-family=true,per-mode=symbol}
\emergencystretch=2em
\allowdisplaybreaks
\newcommand{\vct}[1]{\bm{#1}}
\newcommand{\mat}[1]{\bm{#1}}
\newcommand{\R}{\mathbb R}
\newcommand{\T}{^{\mathsf T}}
\newcommand{\Herm}{^{\mathsf H}}
\newcommand{\E}{\mathbb E}
\newcommand{\Rea}{\operatorname{Re}}
\newcommand{\Ima}{\operatorname{Im}}
\newcommand{\InputPart}[1]{\input{#1}}

\title{A 27-State Directionally Coupled CARMA Error-State Kalman Filter for
Marine Motion: Model, Startup, Ablations, and Degradation}
\author{\IEEEauthorblockN{Mikhail S. Grushinskiy}
\IEEEauthorblockA{Independent Researcher, Bareboat Necessities Project, USA\\
mgrouch@users.github.com}}

\begin{document}
\maketitle

\begin{abstract}
Independent-axis stochastic wave models discard the phase-coherent
horizontal--vertical response of a vessel to directional waves. We study an
experimental quaternion multiplicative extended Kalman filter (QMEKF) with
a six-state real CARMA motion realization on each axis, three attitude-error
states, and six sensor-bias states: 27 error states altogether. The
realization retains the cross-spectral law of a 33-state complexified
shared-mode filter while removing six deterministically unobservable
coordinates. Its lossless two-state colouring cascade has an additional
property: when cross-axis coupling is disabled, its \emph{adaptive} model
reduces exactly in law to the 21-state independent-axis CARMA(4,1) baseline.
We document the actual initialization, excitation-confidence logic and
covariance propagation; then use paired, seeded vessel-response simulations
to study both normal performance and failures. Across 26 seas and 30 seeds
per sea, the 27-state cascade improves three-dimensional RMS error over the
33-state filter in 20 cases, is effectively tied in four and loses in two.
The shared proxy yields the same Live-entry time for all filters, but
startup heave error remains markedly larger than steady-state error.
Coloured accelerometer disturbances and heave-response mismatch reveal
important limitations shared by the realizations. A model of the embedded
arithmetic gives 42\% fewer prediction/update multiply-adds and 21\% less
filter-object memory than the 33-state reference. These are simulation and
compiled-code results, not a claim of sea-trial accuracy or stability proof.
\end{abstract}

\begin{IEEEkeywords}
CARMA processes, directional seas, embedded inertial navigation, marine
heave, quaternion error-state Kalman filter, stochastic realization.
\end{IEEEkeywords}

\section{Introduction}
\label{sec:intro}
An inertial unit on a vessel measures rotational motion and the sum of
gravity, wave-induced acceleration and sensor error. Two integrations
recover relative displacement only if low-frequency acceleration errors
are controlled. Treating surge, sway and heave as three independent scalar
processes provides such regularization but ignores coherence and phase
between motions excited by the same travelling waves. In a directional
sea, the response-amplitude operator (RAO) of a hull generates correlated
translation even when wave components have random phase \cite{jonswap73}.

The ocean-IMU CARMA family uses the same quaternion error-state update,
sensor-bias estimation, magnetic aiding and measurement-only sea-state
tuner as its OU--III family. Its shared-mode extension represents directional
coupling through a Hermitian three-axis driving covariance. The initial
complexified model uses eight real wave states per axis and 33 error states
in total. A prior minimal-realization investigation establishes that two
states per axis are unobservable and compares quotient and cascade
coordinates \cite{kalman63,linpicci15}. The present article makes the
\emph{implemented 27-state filter} the object of a complete system study:
what is propagated, how it starts, which comparisons isolate state reduction
from other changes, and where it degrades.

Three distinctions guide the evaluation. First, a fixed-parameter
stochastic equivalence does not guarantee identical trajectories under
time-varying coefficients. Second, common startup and sensor-front-end
errors cannot be credited to a smaller state realization. Third, an
improvement relative to the 33-state filter may be negligible next to a
failure shared by both; absolute errors are therefore reported alongside
paired differences.

\section{Model and Minimal State Representation}
\label{sec:model}
\subsection{Physical motion and driving spectrum}
Let $\vct p(t)\in\R^3$ denote vessel translation relative to a local wave
equilibrium, $\vct v=\dot{\vct p}$ and $\vct a=\ddot{\vct p}$, expressed in
the filter's world frame. It is not an Earth-fixed navigation position.
The prototype independent-axis CARMA(4,1) transfer from a white-noise input
to each displacement is
\begin{equation}
 H_4(s)=\frac{B(s)}{P(s)},\qquad B(s)=s+\omega_{z0},\quad
 P(s)=P_+(s)\overline{P_+(\bar s)},
 \label{eq:H4}
\end{equation}
where the stable quartic $P$ and the numerator zero are selected by the
C41J shape and scaled by the measured operating frequency. In the C41J
configuration $\omega_{z0}=0$. The 4-state real kinematic realization
$\vct z_a=[p_a,v_a,a_a,\ell_a]\T$ satisfies
$\mathrm d\vct z_a=\mat A_4\vct z_a\,\mathrm dt+
\vct e_3\,\mathrm d\eta_a$ with transfer $H_4$.
Here $\ell_a$ is a CARMA latent coordinate; it is not the OU--III
displacement integral $S_a=\int p_a\,\mathrm dt$.

Directional dependence enters through three circular complex driving
increments $\mathrm d\xi_a=\mathrm d\xi_{r,a}+i\,\mathrm d\xi_{i,a}$,
whose joint spectral covariance is a positive-semidefinite Hermitian
matrix $\mat\Sigma=(\Sigma_{ab})\in\mathbb C^{3\times3}$. The real and
quadrature cross-terms describe co- and phase-shifted axis responses.
The original eight-state realification per axis has transfer
\begin{equation}
 \zeta_a(s)=\frac{B(s)}{P_+(s)D(s)}\xi_a(s),\qquad
 p_a=\Rea\zeta_a,
 \label{eq:complex}
\end{equation}
where $D$ is the stable real quadratic satisfying
$|D(i\omega)|^2=\bigl(|P_+(i\omega)|^2+
|\overline{P_+(-i\omega)}|^2\bigr)/2$.
This expression specifies the model transfer, not an additional measured
wave state. The factorization is a specialized stochastic-realization
construction; related multivariate CARMA factorization is discussed in
\cite{fasen21}.

\subsection{Why six states per axis suffice}
Write $P_+(s)=N_r(s)+iN_i(s)$ for real polynomials $N_r,N_i$. Then
$P=N_r^2+N_i^2$ and the observed real output of \eqref{eq:complex} is
\begin{equation}
 p_a(s)=\frac{B(s)}{P(s)}\eta_a(s),\qquad
 \eta_a(s)=\frac{[N_r(s)\ \ N_i(s)]}{D(s)}
 \begin{bmatrix}\xi_{r,a}(s)\\\xi_{i,a}(s)\end{bmatrix}.
 \label{eq:factor}
\end{equation}
Thus four CARMA kinematic states require only two additional colouring
states. The real quadratic $D$ contributes two repeated eigenvalue
directions to the eight-state realification; the scalar observation cannot
distinguish one direction in each repeated eigenspace. The corresponding
two-dimensional invariant nullspace per axis is independent of $\Sigma$
and of attitude. For fixed model parameters its removal preserves all
observable estimates, innovations, and their covariance, including
three-axis magnetic/acceleration updates. A change in frequency scale
moves this nullspace, so that statement must \emph{not} be extended without
qualification to arbitrary adaptation sequences.

The implemented cascade uses the \emph{co-inner} colouring factor
$\mat V(s)=[N_r(s),N_i(s)]/D(s)$, for which
$\mat V(i\omega)\mat V(i\omega)\Herm=1$. Its controllability-normalized
two-state realization $(\mat A_w,\mat B_w,\mat C_w,
[1,0])$ obeys
\begin{equation}
 \mat A_w+\mat A_w\T=-\mat B_w\mat B_w\T,\qquad
 \mat C_w=-\vct e_1\T\mat B_w\T.
 \label{eq:inner}
\end{equation}
The six-state axis block $\vct x_a=[p_a,v_a,a_a,\ell_a,w_{a1},w_{a2}]\T$
has
\begin{equation}
 \mat A_6=\begin{bmatrix}
 \mat A_4 & \vct e_3\mat C_w\\
 \mat 0 & \mat A_w
 \end{bmatrix},\qquad
 \mat G_6=\begin{bmatrix}
 \vct e_3\vct e_1\T\\ \mat B_w
 \end{bmatrix}.
 \label{eq:A6}
\end{equation}
The full \emph{error} state (the nominal attitude is a separate unit
quaternion) is
\begin{equation}
 \delta\vct x=\left[\delta\vct\theta\T,\ \delta\vct b_g\T,\
 \vct x_x\T,\ \vct x_y\T,\ \vct x_z\T,\
 \delta\vct b_a\T\right]\T\in\R^{27}.
 \label{eq:states}
\end{equation}
The 33-state counterpart replaces each $\vct x_a$ by eight real states;
the 21-state A41 uses only the four-state $\vct z_a$ on each axis.
Twenty-seven is the fixed-coordinate minimal realization preserving the
same six-state \emph{per-axis} shape for any Hermitian excitation. It is
not the minimal stochastic order for one particular excitation: that
order is $4\operatorname{rank}\Sigma+
2\operatorname{rank}\Ima\Sigma$ for the motion law, at most 16, or 25
with the nine attitude/bias coordinates. The latter would require
excitation-dependent coordinates that become singular at zero quadrature
coupling and move as coupling confidence changes.

\subsection{Correlated process covariance}
Let $\mat J_2=\left[\begin{smallmatrix}0&-1\\1&0\end{smallmatrix}\right]$.
For axes $a,b$, the continuous two-input driving cross-covariance is
$g^2(\Rea\Sigma_{ab}\mat I_2+
\Ima\Sigma_{ab}\mat J_2)\mathrm dt$.
The exact discrete model at the IMU interval $h$ is
\begin{align}
 \mat\Phi_6(h)&=e^{\mat A_6h},\\
 \mat Q_{ab}(h)&=\Rea\Sigma_{ab}\mat Q_R(h)
                   +\Ima\Sigma_{ab}\mat Q_I(h), \label{eq:Qab}\\
 \mat Q_{R/I}(h)&=\int_0^h e^{\mat A_6t}\mat W_{R/I}
                                    e^{\mat A_6\Tt}\,\mathrm dt,\\
 \mat W_R&=g^2\mat G_6\mat G_6\T,\quad
 \mat W_I=g^2\mat G_6\mat J_2\mat G_6\T.
\end{align}
$\mat Q_R$ is symmetric and $\mat Q_I$ skew-symmetric. Consequently
$\mat Q_{ba}=\mat Q_{ab}\T$; the assembled 18-state motion covariance
is symmetric positive semidefinite when $\Sigma\succeq0$. Omitting
$\mat Q_I$ would discard the quadrature coupling whose preservation
motivates the model. The implementation discretizes both integrals by a
paired scaled-series/doubling method related to exact linear-SDE
discretization \cite{vanloan78}; IMU and bias process-noise blocks complete
the 27-state covariance.

\section{Estimator, Tuning and Directional Confidence}
\label{sec:method}
\subsection{Quaternion error-state update}
The nominal quaternion $q_{bw}$ rotates body coordinates into the filter
world frame. With angular-rate and accelerometer biases $\vct b_g,\vct b_a$,
the idealized measurements have the form
\begin{equation}
 \vct\omega_m=\vct\omega+\vct b_g+\vct n_g,\quad
 \vct f_m=R(q_{bw})\T(\vct a-\vct g)+\vct b_a+\vct n_a,
 \label{eq:imu}
\end{equation}
and the magnetometer measures a rotated local field with sensor noise.
The filter propagates attitude from bias-corrected angular rate and the
state covariance through the continuous motion model. Accelerometer and
magnetometer measurements update the QMEKF; the acceleration update
couples attitude, bias and the motion block. Importantly, the actual
shared-mode instantiation sets \texttt{SH=SZ=false}: it has no $S$ states
and therefore \emph{does not} perform the periodic $S=0$
pseudo-measurement of OU--III. Its displacement regularization is the
stable CARMA stochastic prior itself. The general core exposes optional
position/velocity pseudo-measurement APIs, but they are not an input of
the evaluated default configuration.
Given a linearized residual $\vct r=\vct y-h(\hat x)$, a typical
measurement uses
\begin{equation}
 \mat K=\mat P^-\mat H\T(\mat H\mat P^-\mat H\T+
 \mat R)^{-1},\quad
 \mat P^+=(\mat I-\mat K\mat H)\mat P^-
 (\mat I-\mat K\mat H)\T+\mat K\mat R\mat K\T.
 \label{eq:joseph}
\end{equation}
The small attitude correction is injected multiplicatively and its
error-state covariance is reset. The observed heading reference, gyro and
accelerometer-bias updates, Joseph covariance step, pseudo-measurement
schedule and reset are the same core machinery as in the CARMA3 baseline;
the 27-state variant changes the motion block, not the sensor-update API.

\subsection{Adaptation is part of the model}
The wrapper estimates a wave-band frequency $f_{\rm op}$ and an
acceleration-evidence statistic from its \emph{measured} leveled IMU data.
The frequency scale $\kappa$ places the CARMA displacement zero-crossing
frequency at the accepted operating frequency; acceleration evidence
modulates the excitation intensity. In Evidence mode, the nonnegative
excess wave-band variance $E$ is weighted by a detection confidence
$p\in[0,1]$, giving an amplitude proportional to
$\sqrt{p\max(E,0)}$. A small numerical floor maintains a nonsingular
stochastic model; it is not interpreted as physical wave excitation.
The motion matrices are normally committed every \SI{0.1}{s}; the IMU
prediction rate is \SI{200}{Hz}. These constants describe the studied
configuration, not a derivation-independent necessity.

When $\kappa$ changes, the cascade holds the physical CARMA coordinates
$[p,v,a,\ell]$ and holds the colouring state after normalization by its
axis noise scale $s_a=\sqrt{g^2}\,\sigma_a$. The corresponding mapping on
the unnormalized coordinates is $w_a\leftarrow
(s_a^{\rm new}/s_a^{\rm old})w_a$. Because \eqref{eq:inner} has a
frequency-independent unit Gramian, the factor remains lossless through
piecewise-constant adaptation. With $\Sigma$ diagonal, this produces
\emph{exactly} the independent-axis A41 stochastic model including
parameter commits. With nonzero quadrature coupling it is a different
time-varying realization than the 33-state complex form, despite their
fixed-parameter equivalence.

\subsection{Coupling estimation, rejection and covariance safety}
Two band cross-spectral regressions estimate
$\vct\beta=[S_{xz}/S_{zz},\ S_{yz}/S_{zz}]\T$ from leveled
acceleration: a long-horizon candidate (nominally \SI{240}{s}) and a
short-horizon consistency monitor (\SI{60}{s}). Coupling is initially
withheld. A confidence $c\in[0,1]$ combines regression agreement,
signal-to-noise, world yaw rate, abrupt frequency change and the
short/long accelerometer innovation ratio. The mismatch term is
\begin{equation}
 D^2=\sum_{j=x,y}\frac{|\beta_{j,\rm cand}-\beta_{j,\rm mon}|^2}
 {v_{j,\rm cand}+v_{j,\rm mon}+\epsilon_j^2},
 \label{eq:confidence}
\end{equation}
with variance estimates $v_j$, a relative floor $\epsilon_j$ and
hysteresis. Asymmetric first-order confidence relaxation uses nominal
time constants of \SI{3}{s} for withdrawal and \SI{60}{s} for recovery.
The confidence scales the candidate coupling at the next model commit.
When it falls, the algorithm also contracts the component of
horizontal--vertical cross-covariance unexplained by common attitude/bias
states. This is a covariance operation, not a reset of position or heading.
Confidence is a safety mechanism for \emph{inconsistent} directional priors,
not a check that the hull RAO used to normalize coupling is correct.

\section{Initialization and Transient Behavior}
\label{sec:startup}
\subsection{Actual startup path}
The 27-state filter is run behind
\texttt{SeaStateFusion\_CarmaSharedModesMinimal}. It inherits the
\texttt{SeaStateFusion\_CARMA3} wrapper and the common
\texttt{ProxyStartupFusion} policy, rather than initializing its wave
states from a known truth trajectory. In Cold/TunerWarm, a measurement-only
Mahony-style proxy \cite{mahony08} estimates tilt and assists magnetic-reference
acquisition while the wave-period tuner collects usable statistics.
Handoff to Live requires the qualified gravity branch, minimum startup
duration and usable tuning and magnetic information as applicable; a
timeout is subject to the gravity-branch gate. The MEKF attitude is seeded
from the proxy, the motion covariance is set to its stationary prior and
accelerometer-bias learning remains gated by accepted magnetic service.
The common implementation also guards excessive tilt and can reinitialize
the attitude after sustained gross misalignment; initial magnetic lock
may be refined later. No special six-state startup algorithm is introduced.

Directional coupling starts at zero confidence, independently of the
Live handoff. Its long candidate/short monitor windows must fill before
the complex cross-spectra may contribute appreciably. It is consequently
incorrect to describe the filter as immediately directionally aided at
Live entry. Equally, the stationarity of the *prior* at initialization is
not evidence that the vessel's first estimated displacement equals the
instantaneous true displacement.

\subsection{Measured startup performance}
The benchmark stores the first Live time $t_L$, the RMS and peak errors in
$[t_L,t_L+120\,\mathrm{s}]$, and a later convergence diagnostic. We define
$E_{z,\rm early}$ as the first-window heave RMS and $E_{z,\rm late}$ as
the RMS after \SI{300}{s}. The convergence statistic uses \SI{30}{s}
rolling RMS: its first reported time is the earliest one after Live from
which \emph{every remaining} rolling window is no worse than
$1.5E_{z,\rm late}$. Its strict whole-future requirement makes it sensitive
to later events and forbids interpretation as a guaranteed physical
settling time.

\begin{table*}[!t]
\caption{Startup in the nine \emph{fixed} seas ($9\times30=270$ paired
realizations per filter). Each entry is a median across runs, except the
95th-percentile Live time and the fraction meeting the strict convergence
criterion within 300 s \emph{after Live}. ``Early/late'' is the median of
the \emph{within-run} RMS ratios, not a ratio of pooled RMS values.}
\label{tab:startup}
\centering\small
\InputPart{csm27-paper-startup.tex-part}
\end{table*}

\begin{figure*}[!t]
\centering
\includesvg[width=0.96\textwidth]{csm27-paper-startup}
\caption{Left: heave RMS during the first 120 s after Live versus scored
RMS on nine fixed seas, each point a 30-seed scenario median (diagonal:
equality). Right: a separately recorded ten-seed ensemble in a
changing-frequency sea, showing 10 s-smoothed, per-second ensemble heave
error. Before Live no QMEKF estimate is plotted.}
\label{fig:startup}
\end{figure*}

The Live time agrees \emph{exactly within each paired run} across all six
filter variants, including the 26-sea evaluation set; the fixed-sea median
is \SI{31.84}{s} (\cref{tab:startup}). The default cascade's early
heave RMS is a median \SI{36.74}{cm} against a median late
\SI{1.93}{cm}; the within-run ratio is 17.3. The large transient is not
removed by six fewer error states. Only 51.9\% of fixed-sea runs pass the
stringent heave convergence criterion within 300 s of Live. This is a
measured limitation of the \emph{common} handoff/low-frequency estimator,
not evidence for a unique defect in the 27-state motion block. A dedicated
initial-tilt/magnetic-dropout startup-policy ablation was not performed
for CSM27 and is not implied by these aggregate numbers.

\section{Evaluation Method}
\label{sec:protocol}
The paired numerical study uses random-phase directional incident waves
and the estimated RAOs of a nominal 28 ft fin-keel vessel at zero forward
speed. Truth motion comes from the wave/RAO simulator, \emph{not} from
the CARMA prior being tested. The same seeded \SI{200}{Hz} IMU and
magnetometer samples feed every estimator in a comparison. The 26 seas
include fixed JONSWAP, Pierson--Moskowitz, swell, regular, crossing and
bimodal cases; wave spreading and headings; gradual and abrupt changes;
and sensor/environmental disturbances. Each has 30 seeds and a nominal
\SI{1200}{s} run. Normal-performance RMS, NEES and innovation statistics
are scored from \SI{300}{s} onward; startup uses the separate window
above. No coefficients of CSM27 are fitted to the evaluation seeds.

The arms are A41@C (21-state CARMA3 baseline with Evidence excitation),
CSM (33-state complex shared modes), CSM27 (cascade), CSM27Q (27-state
quotient with stationary regression lift at frequency commits), CSM27H
(held quotient, ablation) and CSM:off (33-state shared mode with coupling
disabled). They use the same measurement and startup front end.
Changes within a sea are evaluated with matched seeds; 95\% intervals in
the ablation plot are percentile bootstraps of the \emph{paired seed rows},
not samples from the \SI{200}{Hz} time series. The error measures are
\begin{equation}
 E_{3\mathrm D}=\left(\frac1N\sum_k
 \|\hat{\vct p}_k-\vct p_k\|_2^2\right)^{1/2},\quad
 \operatorname{NEES}_p=\vct e_p\T\mat P_{pp}^{-1}\vct e_p.
 \label{eq:metrics}
\end{equation}
The study's per-axis RMS aggregation gives an equivalent $E_{3\mathrm D}$
when all axes use the same scored samples. The nominal NEES reference is
three for a well-specified 3-D Gaussian posterior; an uncalibrated or
time-correlated residual invalidates naive independent-sample confidence
statements. Figures and tables in this article are derived from the
committed study CSVs, which have a revision manifest, rather than from
freshly simulated or inferred data.

\section{Numerical Results and Ablations}
\label{sec:results}
\subsection{Performance across seas}
For the 26-sea comparison, the cascade's paired 3-D RMS is lower than the
33-state filter's on 20 seas by approximately 0.1--0.9\%, statistically
unresolved on four, and higher on two: halved heave RAO and coloured
accelerometer noise. The quotient with a regression lift is within
0.12\% of CSM in 3-D RMS on every sea and within 0.02\% on 24,
empirically isolating the negligible information lost in the removed
coordinates during frequency adaptation. CSM27H, which simply holds
the quotient coordinates at a commit, can lose as much as 1.25\%.
These results concern parameter-adaptive trajectories; they do not
contradict fixed-parameter equivalence.

\begin{table}[!t]
\caption{Mean of scenario-level paired percentage improvements, relative
to CSM (33 states). Positive means \emph{smaller} RMS; there are 9 fixed,
6 coupling, 1 gradual, 3 rapid and 7 disturbed seas.}
\label{tab:groups}
\centering\footnotesize
\resizebox{\columnwidth}{!}{\InputPart{csm27-groups.tex-part}}
\end{table}

\begin{figure*}[!t]
\centering
\includesvg[width=0.94\textwidth]{csm27-paper-ablation}
\caption{Matched-seed realization ablation on seven representative seas:
percentage reduction $100(1-E_f/E_{33})$ of mean 3-D RMS, with 95\%
paired-bootstrap intervals. Positive is better; no point on this plot
measures the absolute severity of a disturbed sea.}
\label{fig:ablation}
\end{figure*}

\subsection{What is gained by the cascade?}
The default cascade is more than an algebraic compression. Its normalized
lossless colouring makes the zero-coupling adaptive law exactly A41,
whereas the 33-state complex coordinates retain an adaptation-specific
``realization effect.'' In spread JONSWAP the cascade improves heave RMS
over CSM by 1.49\% (paired 95\% interval 1.36--1.62\%); on a regular
wave the heave gain is 3.33\%. However, horizontal RMS is slightly worse
on many JONSWAP-like seas. Those changes must be attributed to which
coordinates the tuner holds across frequency commits, not to an
approximation in the fixed continuous-time transfer function.

Three experiments separate mechanisms. First, CSM27Q retains the
observable quotient with a stationary-regression lift at changes of
$\kappa$, closely replicating CSM while using 27 states. Second, CSM27H
removes that lift and tests the assumption that quotient coordinates may
simply be held. Third, a separate seven-sea ablation holds the \emph{raw}
cascade colouring state instead of its noise-normalized value. Raw hold
does not remove the halved-RAO loss (3-D change $-1.84\%$ versus
$-2.15\%$ for the normalized hold) or the coloured-noise loss
($-1.83\%$ versus $-1.61\%$). The default normalized hold is retained
because it preserves the independent-axis fallback equivalence. None of
these changes redefines the shared proxy startup.

\subsection{Uncertainty and independent replay}
\begin{table}[!t]
\caption{Pooled consistency diagnostics over 26 seas and 30 seeds. Values
are medians of per-run time-averaged NEES and NIS, mean empirical heave
95\% interval coverage, and number of failed runs.}
\label{tab:consistency}
\centering\footnotesize
\resizebox{\columnwidth}{!}{\InputPart{csm27-consistency.tex-part}}
\end{table}

The 27-state cascade's median position NEES is 1.652 versus 1.670 for CSM,
and its heave 95\% interval contains truth in 91.2\% of scored samples
(\cref{tab:consistency}). The small median accelerometer NIS (0.091)
also indicates that nominal chi-square expectations do not calibrate these
measurements directly; reporting it is preferable to equating absence of
numerical failure with statistical optimality. No filter failed in the
4680-row study. Eight separately pinned wave/RAO replay records also pass
CSM27's regression limits. These are simulation and regression gates, not
at-sea validation. In the records, the maximum heave error is 2.455\% of
incident $H_s$ for JONSWAP and 2.403\% for PM-Stokes, compared with
2.518\% and 2.458\% for CSM, respectively. The significant gain over
the legacy-excitation TFG-CARMA3 heave score is largely attributable to
\emph{Evidence excitation}, not necessarily to the 27-state realization.

\section{Degradation and Failure Modes}
\label{sec:degrade}
\begin{table*}[!t]
\caption{Absolute 3-D position RMS from matched 30-seed simulations.
Positive values in the last two columns mean \emph{worse} CSM27 error:
$100(E_{27}/E_{\rm comparator}-1)$. This sign is intentionally opposite
to the improvement measure in \cref{fig:ablation}. Disturbed scenarios
alter the data-generating process; they are not controlled one-factor
parameter variations of an identical truth trajectory.}
\label{tab:stress}
\centering\small
\InputPart{csm27-paper-degradation.tex-part}
\end{table*}

\begin{figure*}[!t]
\centering
\includesvg[width=0.96\textwidth]{csm27-paper-degradation}
\caption{Absolute error in eight representative seas, plotted on a log
scale to reveal common-mode failures. The disturbed-scenario error cannot
be inferred from the tiny differences between 27 and 33 states alone.}
\label{fig:stress}
\end{figure*}

\subsection{Noise, motion and measurement disturbances}
The nominal spread JONSWAP case has mean CSM27 3-D RMS of 25.4 mm.
An imposed coloured acceleration disturbance of $0.004g_0$ (approximately
$0.039\,\mathrm{m\,s^{-2}}$), with \SI{2}{s} correlation, changes that
result to 817.5 mm, while A41 and CSM
also rise to 818.6 and 804.6 mm, respectively
(\cref{tab:stress,fig:stress}). This roughly 32-fold CSM27 increase is a
\emph{shared model/measurement weakness}, not an effect of the six removed
states. The paired experiment does not identify a unique causal component;
a noise-floor/colour-model intervention would be needed to do so.

Manoeuvring produces 114.8 mm on CSM27, compared with 115.5 mm on A41,
and an abrupt change in frequency/direction/height produces 64.2 mm on
CSM27. By contrast, $\pm25\%$ sample-interval jitter, magnetic
disturbances with calibration residuals and bias drift leave the observed
3-D errors near 25--27 mm in their tested settings. These finite stress
cells do not establish robustness to arbitrary manoeuvres, jitter, magnetic
outages or calibration error. They identify specific cases in which the
adaptive 27-state prior is and is not the dominant limitation.

\subsection{Wrong directional prior and loss of coherence}
A separate 20-seed robustness grid on three seas supplies deliberate
\emph{external} coupling priors: reversed direction, $90^\circ$ rotated,
$45^\circ$ phase-shifted and doubled coupling amplitude. With the
confidence monitor enabled, both shared-mode filters withdraw inconsistent
coupling and return near the A41 independent-axis error. With confidence
disabled, the doubled prior yields about 57 mm 3-D RMS and position NEES
near 14 in the nominal sea, against about 25 mm with correct coupling.
This ablation demonstrates the value of the confidence state machine, not
an unconditional guarantee: a confidently wrong prior that also agrees
with the short monitor could still pass its gates. The study does not
cover every possible magnetic outage or coherent machinery vibration.

\subsection{A structural heave--surge failure}
With the benchmark hull's heave RAO scaled by 0.5, CSM27 has mean 3-D
error 32.5 mm versus A41's 20.6 mm: a \emph{58\% excess} relative to A41
(\cref{tab:stress}). The same defect affects the 33-state shared mode,
so removing unobservable directions is not the cause. The current
excitation mapping preserves the measured band regression
$\vct\beta=\vct S_{hz}/S_{zz}$ but combines it with independently
normalized horizontal/vertical priors. Its applied coherence then scales
approximately as
\begin{equation}
 \gamma_{\rm applied}^2\propto
 \gamma_{\rm measured}^2
 \frac{S_{hh}/S_{zz}}{\sigma_h^2/\sigma_z^2},
 \label{eq:overcouple}
\end{equation}
so a changed heave-to-horizontal RAO ratio can cause the applied prior to
over-couple even when the monitor observes self-consistent regression.
In the recorded test, measured multiple coherence remains about 0.87,
while applied coherence rises from 0.13 to 0.38 on halving heave. A
measured-coherence cap is a falsifiable remedy but \emph{has not been
implemented or validated here}. This failure would be hidden by comparing
only CSM27 with CSM, whose errors differ by roughly two percent.

\section{Embedded Cost, Scope and Discussion}
\label{sec:embedded}
\begin{table}[!t]
\caption{Structured analytic arithmetic for the model prediction and
accelerometer update, and for a \SI{0.1}{s} model commit. The final column
amortizes commits over the \SI{200}{Hz} sample rate.}
\label{tab:cost}
\centering\footnotesize
\resizebox{\columnwidth}{!}{\InputPart{csm27-cost.tex-part}}
\end{table}

The dense covariance drops from $33^2=1089$ to $27^2=729$ floats
(4356 to 2916 bytes). The structured operation count falls from 13,686
to 7,992 multiply-adds per sample, a 42\% saving; per-commit cost
falls from 28,672 to 19,440 (32\%). The measured ESP32-S3 compilation
with Espressif GCC (\texttt{-Os}) gives a 15,776-byte CSM filter object
against 12,528 bytes for CSM27, a 21\% saving. The trade-off is code
growth from 196,266 to 206,920 bytes (+5.4\%) for the pair-discretization
and closed-form two-state construction. The worst static update call chain
is comparable (4896 versus 4864 bytes). Those are compile-time memory
and instruction-count results, \emph{not} measured ESP32-S3 deadlines,
power consumption or endurance. Host timing was sensitive to denormal
floating-point handling, a reason not to substitute x86 latency for MCU
latency.

There are important scope restrictions. The vessel response is an
estimated synthetic RAO, the wave realization is random-phase simulation,
the test magnetometer and calibration errors are synthetic, and the
reference records do not measure free-running real-vessel truth. The
startup comparison is observational across a shared policy, rather than
a controlled feature-off experiment. The published confidence bounds
condition on the chosen scenarios and seeded simulator; they do not
measure model uncertainty over unobserved hulls or sensor installations.
No formal stability theorem is claimed for the adaptive, nonlinear,
gated, finite-precision implementation. Only the OU--III filter is the
project's shipping default; CSM27 remains experimental.

\section{Conclusion}
The 27-state shared-mode CARMA QMEKF preserves the fixed-parameter
three-axis stochastic law of the 33-state complexified implementation,
including quadrature covariance, and substantially reduces its embedded
covariance work. The cascade is a practical default among the 27-state
realizations because its zero-coupling adaptive law coincides with the
independent-axis CARMA3 baseline and its model commits are well conditioned.
The transient and stress experiments impose a more qualified conclusion:
startup remains a common front-end limitation; a coloured low-frequency
acceleration disturbance defeats all three CARMA variants tested; and
the present coupling normalization can harm estimation on a hull with
unexpected heave response. Further work should pair a measured-coherence
cap with a full 26-sea re-test, isolate the colored-noise excitation
mechanism, test initial-attitude/magnetometer startup ablations, and
measure timing and behavior on instrumented vessels.

\section*{Reproducibility and Provenance}
The implementation is in
\texttt{src/kalman\_carma\_shared\_modes/CarmaSharedModesMinimalMotion.h},
\texttt{KalmanCarmaSharedModesMinimal.h} and
\texttt{SeaStateFusionFilter\_CarmaSharedModesMinimal.h}.
The full experiment driver is
\texttt{tools/carma\_shared\_modes\_minimal\_study.py} with
\texttt{tests/kalman\_carma\_shared\_modes/}. The committed study
\texttt{reports/results/kalman\_carma\_shared\_modes\_minimal/}
contains the 4680-run primary CSV, 630-run commit-rule ablation,
1240-run robustness grid, replays, trace envelopes, hardware size probe
and a manifest identifying the original source revision. This article's
\texttt{csm27-paper-*.tex-part} tables and vector PDF figures are
reproducibly regenerated by
\texttt{python3 tools/csm27\_system\_paper.py} from those committed
files. The other shared-mode papers in this directory give the full
observability and state-factorization proofs and the preceding 33-state
validation; the present article does not present those derivations as
new experimental results. Project repository:
\url{https://github.com/bareboat-necessities/ocean-imu}.

\begin{thebibliography}{9}
\bibitem{kalman63}
R. E. Kalman, ``Mathematical description of linear dynamical systems,''
\emph{J. SIAM Control}, vol. 1, no. 2, pp. 152--192, 1963.
\bibitem{linpicci15}
A. Lindquist and G. Picci, \emph{Linear Stochastic Systems: A Geometric
Approach to Modeling, Estimation and Identification}. Springer, 2015.
\bibitem{vanloan78}
C. F. Van Loan, ``Computing integrals involving the matrix exponential,''
\emph{IEEE Trans. Autom. Control}, vol. 23, no. 3, pp. 395--404, 1978.
\bibitem{jonswap73}
K. Hasselmann \emph{et al.}, ``Measurements of wind-wave growth and swell
decay during the Joint North Sea Wave Project (JONSWAP),''
\emph{Deutsche Hydrographische Zeitschrift, Supplement A}, no. 8, 1973.
\bibitem{mahony08}
R. Mahony, T. Hamel, and J.-M. Pflimlin, ``Nonlinear complementary filters
on the special orthogonal group,'' \emph{IEEE Trans. Autom. Control},
vol. 53, no. 5, pp. 1203--1218, 2008.
\bibitem{fasen21}
V. Fasen-Hartmann and M. Scholz, ``Factorization and discrete-time
representation of multivariate CARMA processes,'' arXiv:2102.11681, 2021.
\end{thebibliography}

\end{document}
